% Введение

\section*{Введение}
\addcontentsline{toc}{section}{Введение}

Актуальность темы исследования определяется стремительным развитием генеративных моделей на основе диффузионного процесса и возросшим практическим интересом к задаче художественного переноса стиля (style transfer). Современные диффузионные модели достигли фотореалистичного качества синтеза изображений, однако управляемое разделение стилистических и семантических характеристик по-прежнему остаётся нетривиальной задачей. Традиционные методы переноса стиля, основанные на свёрточных нейронных сетях~\cite{gatys2016nst}, не позволяют достичь необходимой гибкости при работе с современными трансформерными архитектурами. Появление диффузионных трансформеров (Diffusion Transformer, DiT)~\cite{peebles2023dit} и многомодальных вариантов архитектуры (Multimodal Diffusion Transformer, MM-DiT) открыло новые возможности для исследования внутреннего пространства представлений генеративных моделей.

Модель \textbf{FLUX.1}~\cite{blackforest2024flux}, включающая порядка 12 миллиардов параметров и построенная на базе архитектуры MM-DiT с применением выпрямленных потоков (rectified flow), представляет собой одну из наиболее мощных открытых генеративных систем на сегодняшний день. Её архитектура включает 19 блоков двупоточной обработки (double-stream blocks) и 38 блоков однопоточной обработки (single-stream blocks), что принципиально отличает её от предшествующих моделей на основе U-Net. Несмотря на превосходство FLUX.1 в качестве синтеза, методы эффективной тонкой настройки (fine-tuning) с разделением стиля и содержания для данной архитектуры остаются недостаточно изученными.

Метод \textbf{B-LoRA}~\cite{frenkel2024blora}, предложенный для модели SDXL, продемонстрировал принципиальную возможность разделения стилистических и контентных характеристик посредством избирательного обучения адаптеров низкого ранга (Low-Rank Adaptation, LoRA) в различных блоках U-Net архитектуры. Перенос данного принципа на трансформерную архитектуру FLUX.1 представляет как теоретический, так и практический интерес: MM-DiT обрабатывает текстовые и визуальные токены принципиально иным образом, нежели U-Net, что требует переосмысления гипотезы о специализации блоков.

\textbf{Цель} настоящей работы --- разработать и экспериментально обосновать метод адаптации принципа B-LoRA к архитектуре FLUX.1 для управляемого разделения стиля и содержания при художественном переносе стиля, превосходящий по качеству существующие подходы при сопоставимых вычислительных затратах.

Для достижения поставленной цели необходимо решить следующие \textbf{задачи}:

\begin{enumerate}
    \item Провести систематический анализ существующих методов художественного переноса стиля в диффузионных моделях, включая подходы на основе LoRA, IP-Adapter~\cite{ye2023ipadapter} и конкурирующий метод SplitFlux~\cite{yang2025splitflux}, и выявить их ограничения применительно к архитектуре FLUX.1.
    \item Сформулировать гипотезу о функциональной специализации блоков модели FLUX.1: соответствие блоков двупоточной обработки (блоки 0--5 как начальная гипотеза, уточняемая в ходе аблационного исследования, эксперимент~A) кодированию контентной информации, а блоков однопоточной обработки (блоки 30--37 как начальная гипотеза, также проверяемая экспериментом~A) --- стилистической.
    \item Разработать метод поблочного обучения адаптеров LoRA в рамках архитектуры FLUX.1 с раздельными наборами параметров $\Theta_{\mathrm{content}}$ и $\Theta_{\mathrm{style}}$, предусматривающий совместное обучение обоих адаптеров --- в отличие от метода SplitFlux~\cite{yang2025splitflux}, ограничивающегося исключительно стилевыми блоками, --- а также схему инференса с использованием исключительно стилевого адаптера.
    \item Сформировать экспериментальную базу: подготовить обучающие наборы данных изображений в стиле Ван Гога и Моне, а также оценочные наборы на основе COCO-100~\cite{lin2014coco} и ArtBench-10~\cite{liao2022artbench}.
    \item Провести аблационное исследование (эксперименты A, B, C) для проверки значимости вклада отдельных компонент предложенного метода, включая систематическую проверку гипотезы о выборе диапазонов блоков.
    \item Выполнить сравнительную оценку предложенного метода относительно базовых подходов --- Full-LoRA-FLUX, IP-Adapter-FLUX и B-LoRA-SDXL --- по метрикам CLIP-style, CLIP-content, FID и LPIPS.
    \item Проанализировать полученные результаты и сформулировать выводы об условиях применимости и ограничениях предложенного метода.
\end{enumerate}

\textbf{Объектом исследования} является диффузионная трансформерная модель FLUX.1 как представитель класса архитектур MM-DiT с выпрямленными потоками.

\textbf{Предметом исследования} служат методы поблочной адаптации низкорангового дообучения (LoRA) в рамках архитектуры FLUX.1, обеспечивающие разделение пространств стиля и содержания при художественном переносе стиля.

\textbf{Методы исследования} включают: анализ архитектурных особенностей диффузионных трансформеров; методы параметрически эффективного дообучения (Parameter-Efficient Fine-Tuning, PEFT); количественную оценку качества генерации с использованием мультимодальных метрик на основе модели CLIP~\cite{radford2021clip}; методы статистического сравнения экспериментальных результатов. Реализация выполнена с использованием фреймворков PyTorch~2.1~\cite{paszke2019pytorch}, \texttt{diffusers}~\cite{von2022diffusers}, \texttt{peft}~\cite{peft2022}, системы управления экспериментами ClearML и инструмента \texttt{ai-toolkit}.

\textbf{Научная новизна} работы состоит в следующем: предложен метод совместного поблочного обучения стилевого и контентного адаптеров LoRA ($\Theta_{\mathrm{style}}$ и $\Theta_{\mathrm{content}}$) применительно к архитектуре FLUX.1/MM-DiT; в отличие от метода SplitFlux~\cite{yang2025splitflux}, в настоящей работе оба адаптера обучаются совместно, а роль блоков двупоточной обработки в кодировании контента проверяется систематически; разработана и теоретически обоснована гипотеза о функциональной специализации блоков двупоточной и однопоточной обработки модели FLUX.1, проверяемая посредством систематического аблационного исследования; предложен подход к переносу принципа B-LoRA с архитектуры U-Net на трансформерную архитектуру с сохранением возможности управляемого разделения стиля и содержания.

\textbf{Практическая значимость} работы определяется тем, что разработанный метод позволяет осуществлять высококачественный художественный перенос стиля на основе модели FLUX.1 при использовании малых обучающих наборов данных без модификации весов базовой модели. Предложенный подход применим в задачах цифрового искусства, персонализированной генерации и промышленного дизайна, а также может быть распространён на другие модели класса MM-DiT.

\textbf{Структура работы.} Работа состоит из введения, четырёх глав, заключения и приложений. Глава~2 содержит обзор литературы по методам художественного переноса стиля и параметрически эффективного дообучения диффузионных моделей. В главе~3 описан предложенный метод поблочного обучения адаптеров LoRA в архитектуре FLUX.1. Глава~4 посвящена экспериментальной проверке метода, аблационному исследованию и анализу полученных результатов. В приложении~А приведено описание использованного технологического стека.
